% ---- p2-namespaced TikZ styles (package loading happens in preamble.txt) ----
\tikzset{
  p2ent/.style={circle, draw=black!65, line width=0.7pt, minimum size=9.5mm,
                inner sep=0pt, font=\small\bfseries},
  p2lab/.style={font=\scriptsize},
  p2rel/.style={draw=black!65, line width=0.7pt},
  p2el/.style={font=\scriptsize, fill=white, inner sep=1.5pt},
}
\newcommand{\pTwoTgt}{orange!22}
\newcommand{\pTwoAgt}{blue!12}
\newcommand{\pTwoAtt}{green!14}
\newcommand{\pTwoVen}{purple!12}

\subsection*{4. Baselines}
\label{p2:sec:baselines}

\subsubsection*{4.1 Homogeneous methods, reported as type-agnostic baselines}

These are reported in essentially every heterogeneous paper, and the \emph{point} of reporting
them is that they \textbf{discard $\phi$ and $\psi$ entirely} --- they run on the graph with all
type information erased (or on one meta-path-induced projection). They are the control condition:
whatever a heterogeneous model gains over them is the measured value of modelling types. HGB's
central and uncomfortable finding is that this gap is much smaller than the literature had
claimed, and that a properly tuned GAT beats most specialised HGNNs.

\begin{itemize}\itemsep2pt
\item \textbf{DeepWalk} --- Perozzi, Al-Rfou \& Skiena, KDD 2014. Truncated uniform random walks
  treated as sentences, fed to skip-gram with hierarchical softmax.
\item \textbf{LINE} --- Tang, Qu, Wang, Zhang, Yan \& Mei, WWW 2015. Explicitly preserves
  first-order (observed edge) and second-order (shared-neighbourhood) proximity, optimised with
  negative sampling and edge-sampling for large graphs.
\item \textbf{node2vec} --- Grover \& Leskovec, KDD 2016. DeepWalk with a biased second-order walk
  whose $p,q$ parameters interpolate between BFS-like (structural equivalence) and DFS-like
  (homophily) exploration.
\item \textbf{GCN} --- Kipf \& Welling, ICLR 2017. First-order spectral convolution: a layer is a
  symmetrically normalised neighbourhood average followed by a linear map and nonlinearity.
\item \textbf{GAT} --- Veli\v{c}kovi\'{c}, Cucurull, Casanova, Romero, Li\`{o} \& Bengio, ICLR
  \textbf{2018}. Replaces the fixed normalisation with multi-head self-attention over neighbours.
  (Commonly miscited as 2017, the arXiv posting date.)
\item \textbf{GraphSAGE} --- Hamilton, Ying \& Leskovec, NIPS 2017. Inductive framework that
  samples a fixed number of neighbours and aggregates with mean/LSTM/max-pool, enabling embeddings
  for unseen nodes.
\end{itemize}

\subsubsection*{4.2 Meta-path and random-walk HIN embedding}

\begin{itemize}\itemsep3pt
\item \textbf{metapath2vec} --- Dong, Chawla \& Swami, KDD 2017. Random walks are
  \emph{meta-path-guided}: at each step the walker must move to the node type prescribed by the
  meta-path scheme (e.g.\ APVPA), so the generated context is type-balanced instead of being
  dominated by the most frequent type. The walks feed a heterogeneous skip-gram. Note that the
  base model still normalises negative sampling over \emph{all} nodes regardless of type; it is the
  \textbf{metapath2vec++} variant that normalises the softmax per node type.
\item \textbf{HIN2Vec} --- Fu, Lee \& Lei, CIKM 2017. Rather than fixing meta-paths in advance, it
  learns node \emph{and} relation (meta-path) vectors jointly by training a multi-task binary
  classifier that predicts, for a node pair, which relations hold between them.
\item \textbf{ESim} --- Shang, Qu, Liu, Kaplan, Han \& Peng, arXiv:1610.09769, 2016
  (\emph{preprint only --- never published in a peer-reviewed venue}). Embeds nodes by sampling
  positive and negative path instances of user-supplied weighted meta-paths and optimising a
  noise-contrastive objective for similarity search.
\item \textbf{HERec} --- Shi, Hu, Zhao \& Yu, IEEE TKDE 31(2):357--370, \textbf{2019}. Performs
  meta-path-guided random walks, but first \emph{filters} each walk to the target type, producing
  homogeneous sequences per meta-path; the resulting embeddings are fused by a nonlinear function
  and injected into an extended matrix-factorisation recommender.
\item \textbf{PTE} --- Tang, Qu \& Mei, KDD 2015. Decomposes a heterogeneous text network
  (word--word, word--document, word--label) into bipartite graphs and jointly embeds them with the
  LINE objective, giving semi-supervised, task-tuned text embeddings.
\end{itemize}

\subsubsection*{4.3 Heterogeneous graph neural networks}

\begin{itemize}\itemsep3pt
\item \textbf{R-GCN} --- Schlichtkrull, Kipf, Bloem, van den Berg, Titov \& Welling, ESWC
  \textbf{2018}. Relation-specific message passing: each layer aggregates neighbours separately per
  relation $r$ with its own weight matrix $W_r$, normalised by $c_{i,r}=|\mathcal{N}_i^r|$, plus a
  self-loop term. Because $|\mathcal{R}|$ weight matrices overfit, it offers two \emph{alternative}
  regularisers --- \emph{basis decomposition} ($W_r$ a linear combination of $B$ shared bases) and
  \emph{block-diagonal decomposition} ($W_r$ a direct sum of small dense blocks). For link
  prediction it is used as an encoder with a \textbf{DistMult} decoder.
\item \textbf{HAN} --- Wang, Ji, Shi, Wang, Ye, Cui \& Yu, WWW 2019. Two-level attention:
  \emph{node-level} GAT-style attention within each meta-path-induced neighbour set, then
  \emph{semantic-level} attention learning a scalar importance per meta-path to fuse them. Its
  limitation --- the one MAGNN attacks --- is that it operates on meta-path-\emph{induced} graphs,
  so intermediate nodes along the meta-path are discarded.
\item \textbf{MAGNN} --- Fu, Zhang, Meng \& King, WWW 2020. Three stages: node-content
  transformation projecting differently-dimensioned type attributes into a shared space;
  \emph{intra}-meta-path aggregation that encodes the \textbf{entire path instance including
  intermediate nodes} (mean, linear, or a relational-rotation encoder borrowed from RotatE) with
  attention over instances; and \emph{inter}-meta-path semantic attention.
\item \textbf{GTN} --- Yun, Jeong, Kim, Kang \& Kim, NeurIPS 2019. \emph{Learns} meta-paths instead
  of accepting them. A $1\times1$ convolution with softmax weights forms a soft convex combination
  $Q_i$ of the per-edge-type adjacency matrices; multiplying $A=Q_1Q_2\cdots Q_l$ yields new
  meta-path adjacencies differentiably. Including the identity among the candidates lets it learn
  paths \emph{shorter} than $l$. A GCN then runs on the learned graph.
\item \textbf{HetGNN} --- Zhang, Song, Huang, Swami \& Chawla, KDD 2019. Restart-based random walks
  sample a fixed-size, type-balanced neighbour set; content of each node is encoded with a
  Bi-LSTM over its heterogeneous attributes, neighbours are aggregated per type, and types are
  combined by attention.
\item \textbf{HGT} --- Hu, Dong, Wang \& Sun, WWW 2020. Node-type-specific projections produce
  Key/Query/Value, and attention and message matrices are indexed by the \emph{meta-relation}
  $\langle\tau(s),\phi(e),\tau(t)\rangle$, with a learnable scalar prior $\mu$ per meta-relation.
  It needs no hand-designed meta-paths. Adds \emph{relative temporal encoding} for dynamic graphs
  and \emph{HGSampling} for web-scale mini-batch training.
\item \textbf{simple-HGN / HGB} --- Lv, Ding, Liu et al., KDD 2021. A deliberately minimal
  baseline: GAT plus (i) a learnable \textbf{edge-type embedding} concatenated into the attention
  logit, (ii) residual connections on \emph{both} node representations and attention scores, and
  (iii) $L_2$ normalisation of the output embeddings. The paper's larger contribution is the HGB
  benchmark itself and the finding that this simple model matches or beats most specialised HGNNs.
\item \textbf{ie-HGCN} --- Yang, Guan, Li, Zhao, Cui \& Wang, IEEE TKDE 35(2):1637--1650, 2023.
  Hierarchical object-level then type-level aggregation which, as a by-product, \emph{ranks} the
  useful meta-paths, making the model interpretable while avoiding an exhaustive meta-path search.
\item \textbf{RSHN} --- Zhu, Zhou, Pan, Zhu \& Wang, IEEE ICDM 2019. Builds a coarsened
  \emph{edge-centric} ``relation structure graph'' to embed relation types themselves, then
  propagates with those relation embeddings rather than treating relations as bare labels.
\item \textbf{HPN} --- Ji, Wang, Shi, Wang \& Yu, IEEE TKDE 35(1):521--532, 2023. Analyses
  semantic confusion (the HIN analogue of over-smoothing) and replaces the propagation step with a
  personalised-PageRank-like update that retains a fraction of the node's own representation,
  enabling deeper meta-path propagation.
\end{itemize}

\subsubsection*{4.4 Self-supervised / contrastive heterogeneous methods}

\begin{itemize}\itemsep3pt
\item \textbf{DMGI} --- Park, Kim, Han \& Yu, AAAI 2020. Extends Deep Graph Infomax to
  \emph{attributed multiplex} networks: one DGI mutual-information objective per relation type,
  with a consensus regularisation that pulls the per-relation embeddings towards a shared
  representation.
\item \textbf{HDGI} --- Ren, Liu, Huang, Dai, Bo \& Zhang, arXiv:1911.08538, 2019 (AAAI-20 DLGMA
  workshop, under a different title; \emph{no archival venue}). Applies mutual-information
  maximisation between local node embeddings and a global summary, using meta-path-based
  convolution plus semantic attention as the encoder.
\item \textbf{HeCo} --- Wang, Liu, Han \& Shi, KDD 2021. Contrasts two \emph{views} of the same
  node: a \emph{network-schema} view (attention over directly-connected neighbours of each other
  type) and a \emph{meta-path} view (meta-path-specific GCNs with semantic attention). Positives
  are not merely the node's counterpart view --- HeCo selects multiple positives by counting how
  many meta-path instances connect two nodes and taking the top ones.
\end{itemize}

\subsubsection*{4.5 Knowledge-graph embeddings as relation-aware LP baselines}

A HIN with many relation types is close to a knowledge graph, so heterogeneous link-prediction
papers (and HGB's LP track) routinely report KG embedding models, which score a typed triple
$(h,r,t)$ directly.

\begin{itemize}\itemsep2pt
\item \textbf{TransE} --- Bordes, Usunier, Garcia-Duran, Weston \& Yakhnenko, NIPS 2013. Models a
  relation as a translation, $\mathbf{h}+\mathbf{r}\approx\mathbf{t}$, scoring
  $-\|\mathbf{h}+\mathbf{r}-\mathbf{t}\|$. Cannot represent 1-to-N or symmetric relations.
\item \textbf{DistMult} --- Yang, Yih, He, Gao \& Deng, ICLR 2015. Bilinear score
  $\langle\mathbf{h},\mathrm{diag}(\mathbf{r}),\mathbf{t}\rangle$; efficient but necessarily
  \emph{symmetric}.
\item \textbf{ComplEx} --- Trouillon, Welbl, Riedel, Gaussier \& Bouchard, ICML 2016. Moves
  DistMult into $\mathbb{C}$ and takes the real part of a Hermitian product, so asymmetric
  relations become representable at the same cost.
\item \textbf{ConvE} --- Dettmers, Minervini, Stenetorp \& Riedel, AAAI 2018. Reshapes $h$ and $r$
  into 2-D and applies a convolution, giving a parameter-efficient, highly expressive scorer.
\item \textbf{RotatE} --- Sun, Deng, Nie \& Tang, ICLR 2019. Treats a relation as an element-wise
  \emph{rotation} in complex space, $\mathbf{t}=\mathbf{h}\circ\mathbf{r}$ with
  $|r_i|=1$, capturing symmetry, antisymmetry, inversion and composition simultaneously.
\end{itemize}

% =====================================================================
